\documentclass[11pt]{article}
\usepackage[margin=1in]{geometry}
\usepackage{amsmath,amssymb}
\usepackage[hidelinks]{hyperref}

\title{Literature Status: Hardy-Space Quasi-Greedy Democracy}
\author{Automatic Math Discovery Run \texttt{fa\_banach\_001}}
\date{June 14, 2026}

\begin{document}
\maketitle

\section*{Original Problem}

The source paper is Fernando Albiac, Jose L. Ansorena, and Przemyslaw
Wojtaszczyk, \emph{Quasi-greedy bases in the spaces \(\ell_p\)
\((0<p<1)\) are democratic}, arXiv:2004.05206; later published in
\emph{J. Funct. Anal.} 280 (2021), no. 7, 108871.

Question 3.8 appears on PDF page 16 of the source paper, after the proof that
quasi-greedy bases in \(\ell_p\), \(0<p<1\), are democratic. It asks:
\[
\text{Let }0<p<1\text{ and }d\in\mathbb N.\text{ Are all quasi-greedy bases in }
H_p(\mathbb D^d)\text{ democratic?}
\]
The same passage is in
\texttt{data/parsed/arxiv\_sources/2004.05206/source.tex}, lines 555--557.

\section*{Supporting Answer Source}

The supporting paper is Fernando Albiac, Jose L. Ansorena, and Glenier Bello,
\emph{Democracy of quasi-greedy bases in \(p\)-Banach spaces with applications
to the efficiency of the TGA in the Hardy spaces \(H_p(\mathbb D^d)\)},
arXiv:2208.09342.

Its abstract says that the authors prove all quasi-greedy bases of
\(H_p(\mathbb D)\), \(0<p<1\), are democratic, while no quasi-greedy basis of
\(H_p(\mathbb D^d)\) for \(d\ge 2\) is democratic, and that this solves a
problem raised in \cite{AAW2021}. In the introduction, PDF page 4, the paper
states explicitly that it solves Question 3.8 from \cite{AAW2021} positively
when \(d=1\), and proves that no quasi-greedy basis of the multivariate Hardy
spaces is democratic when \(d\ge2\). Its bibliography identifies
\cite{AAW2021} as the JFA version of arXiv:2004.05206.

\section*{Answer}

Question 3.8 from arXiv:2004.05206 is already answered in the later literature.
The answer depends on the dimension:
\[
d=1:\quad \text{yes, all quasi-greedy bases of }H_p(\mathbb D)\text{ are democratic;}
\]
\[
d\ge2:\quad \text{no, no quasi-greedy basis of }H_p(\mathbb D^d)\text{ is democratic.}
\]

For \(d=1\), Corollary 4.3 of arXiv:2208.09342, PDF page 21, proves more
generally that if \(\mathcal Y\) is a truncation quasi-greedy basis of a
complemented subspace of \(H_p(\mathbb D)\), then \(\mathcal Y\) is democratic
with upper democracy function equivalent to \(m^{1/p}\). In particular, all
quasi-greedy bases of \(H_p(\mathbb D)\) are almost greedy, hence democratic.

For \(d\ge2\), Corollary 4.6 of arXiv:2208.09342, PDF page 22, proves that
\(H_p(\mathbb D^d)\) has no squeeze-symmetric bases and in particular no
almost greedy bases. Together with the quasi-Banach greedy-basis
characterization used by the authors, this gives the negative answer stated
explicitly in their abstract and introduction: no quasi-greedy basis of
\(H_p(\mathbb D^d)\), \(d\ge2\), is democratic.

\section*{Verification Notes}

Local source text checks:
\begin{itemize}
\item \texttt{data/parsed/arxiv\_sources/2004.05206/source.tex}, lines
555--557, contains the source Question 3.8.
\item \texttt{data/parsed/arxiv\_sources/2208.09342/source.tex}, line 114,
states in the abstract that the paper solves the problem raised in
\cite{AAW2021}.
\item \texttt{data/parsed/arxiv\_sources/2208.09342/source.tex}, line 142,
identifies the source question as Question 3.8 and states the \(d=1\) and
\(d\ge2\) outcomes.
\item \texttt{data/parsed/arxiv\_sources/2208.09342/source.tex}, lines
726--733, contains Corollary 4.3 for \(H_p(\mathbb D)\).
\item \texttt{data/parsed/arxiv\_sources/2208.09342/source.tex}, lines
779--784, contains Corollary 4.6 for \(H_p(\mathbb D^d)\), \(d\ge2\).
\item \texttt{data/parsed/arxiv\_sources/2208.09342/source.tex}, lines
944--956, identifies \cite{AAW2021} as the JFA paper corresponding to
arXiv:2004.05206.
\end{itemize}

\section*{Search Bounds}

Before making this packet, the run registry, solution index, attempt index,
and proof-gap index were searched for arXiv:2004.05206, arXiv:2208.09342,
``Hardy'', \(H_p\), \(H_p(\mathbb D^d)\), and quasi-greedy democracy terms.
No existing packet or ledger row for this exact literature answer was found.
The current target pool already contained arXiv:2208.09342 as a related later
paper. A web search on June 14, 2026 found the arXiv record for
arXiv:2208.09342, whose abstract gives the same explicit resolution.

\section*{Limitations}

This packet records an already-known literature answer, not an original proof
or counterexample produced by the run. It does not claim any resolution of
new questions raised in arXiv:2208.09342.

\section*{Human Review Recommendation}

Check the bibliographic identification of \cite{AAW2021} with
arXiv:2004.05206 and confirm that the source Question 3.8 is the same
Hardy-space democracy question answered in arXiv:2208.09342. The supporting
paper explicitly identifies the problem, so review should be straightforward.

\begin{thebibliography}{9}
\bibitem{AAW2021}
F. Albiac, J. L. Ansorena, and P. Wojtaszczyk,
\emph{Quasi-greedy bases in the spaces \(\ell_p\) \((0<p<1)\) are democratic},
J. Funct. Anal. 280 (2021), no. 7, 108871; arXiv:2004.05206.

\bibitem{AAB2022}
F. Albiac, J. L. Ansorena, and G. Bello,
\emph{Democracy of quasi-greedy bases in \(p\)-Banach spaces with applications
to the efficiency of the TGA in the Hardy spaces \(H_p(\mathbb D^d)\)},
arXiv:2208.09342.
\end{thebibliography}

\end{document}
